% \section{Personas in Language Models}\label{sec:personas-in-language-models}

% \textbf{Persona in the LLM literature.} In this paper, we define a persona as a recurring pattern in model behavior, summarized by the model’s position along a set of measurable trait dimensions. More concretely, let a model with weights $W$ respond to prompts $x$ under context $c$, producing responses from $p_W(y|x,c)$. Given a collection of trait scorers $s_1, \dots, s_k$, we can represent the model’s behavior over a prompt/context distribution as a vector of expected trait scores. A persona is then a region, rather than a point, in this trait space: predictive over some range of contexts and shaped by the immediate prompt, system message, dialogue history, and sampling procedure.

% \textbf{LLMs are trained on substantial amounts of human data.} What might such regions look like? We begin with human personality traits, as human psychometrics provides a well-developed language for measuring stable behavioral variation. Since LLMs are trained on large amounts of human-generated text, they are likely to exhibit many of the regularities associated with human personality descriptions. For example, LLMs are able to act as predictors for human personality traits when analysing the behaviour of characters in stories \citep{suh2024rediscovering}, exhibit bias when prompted to be anxious \citep{codaforno2024inducinganxietylargelanguage}, and react to representations of emotion concepts in patterns consistent with humans under the influence of an emotion \citep{sofroniew2026twheemotion}. 
